\chapter{绪论}

\section{研究背景}

随着我国高速公路网络规模持续扩大，道路交通运行正呈现出高速度、高密度、长距离、跨区域和全天候等特点。国家统计局发布的《中华人民共和国 2025 年国民经济和社会发展统计公报》显示，2025 年末全国民用汽车保有量已达到 36611 万辆，其中私人汽车保有量 32336 万辆；若回看上一年度，2024 年末全国民用汽车保有量为 35268 万辆，新能源汽车保有量已达 3140 万辆。这说明我国机动车规模仍处于高位增长阶段，道路交通系统面临的安全运行压力持续加大。与此同时，交通运输部发布的《2024 年交通运输行业发展统计公报》指出，截至 2024 年底，全国公路总里程达到 549.5 万公里，高速公路里程达到 19.07 万公里。大规模、高等级的路网扩张显著提升了人流、物流和区域经济联系效率，也意味着极端天气或低能见度事件一旦发生，其影响范围、连锁效应和管理难度都将更加突出。

在道路交通安全诸多影响因素中，低能见度天气一直是事故诱因中的高风险因素。世界卫生组织在道路交通伤害事实资料中指出，全球每年约有 119 万人死于道路交通事故。对高速公路而言，影响安全的不仅是降雨、冰雪、横风和路面结冰，雾特别是团雾所引起的视觉退化同样会显著增加风险。美国联邦公路管理局（FHWA）长期研究表明，美国每年有 38700 余起交通事故与雾有关，并造成 600 余人死亡、16300 余人受伤。虽然这一统计并不能直接映射到我国，但它足以说明：在机动化水平较高、交通组织复杂、道路运行密集的国家，雾相关交通风险具有普遍性与高后果性。我国高速公路出行中，自驾与货运比重持续上升，运输时段不断向夜间、清晨等低光照时段延伸，这进一步放大了雾天低能见度带来的安全问题。

团雾是一种在局地地形、地表热力和气象条件共同作用下形成的低能见度天气现象，其最大特征在于“局地发生、强度不均、突发性强、边界模糊”。在同一条高速公路路段上，几十米外可能仍能保持较高能见度，而局部区域已出现严重的视觉遮挡。传统气象监测体系通常更擅长给出宏观区域的天气判断，而对高速公路这种线性基础设施上的局地、快速变化天气识别难度较大。另一方面，现有监控中心大量依赖人工观察视频监控画面或基于规则的阈值策略进行判断，这在复杂夜间背景、车灯干扰、逆光、雾区局地变化等情况下容易产生误判、漏判和延迟处置。

从业务需求来看，高速公路团雾监测至少应满足以下几方面能力：第一，能够在固定视频监控场景下识别道路当前是否处于无雾、均匀雾还是局地团雾状态；第二，能够给出雾强度的相对量化描述，为分级预警、动态限速、车道管控和诱导发布提供依据；第三，能够同时感知车辆目标，从而使系统不仅知道“天气恶化了”，还知道“恶化天气下车辆运行的可视条件与目标识别条件如何变化”；第四，算法必须具有较强工程可落地性，既要考虑训练数据获取成本，也要兼顾推理时的速度、稳定性与部署复杂度。

因此，围绕固定监控视频中的团雾监测构建一个兼顾检测、天气识别与雾浓度估计的统一系统，既具有明显的理论研究意义，也具有较强的实际应用价值。对于管理部门而言，这样的系统可以作为交通气象监测体系的视觉补充；对于道路运营方而言，它有助于实现对高风险雾区路段的精细感知；对于学术研究而言，它则提供了一个将目标检测、低层退化建模、单目深度估计和多任务学习有机结合的问题场景。

\subsection{陕西省交通运输发展现状与课题应用场景}

若进一步从区域落地角度审视本文选题，陕西省具有较强的代表性与典型性。2026 年 1 月 26 日召开的陕西省交通运输工作会议通报显示，截至 2025 年底，全省高速公路通车里程已达 6895 公里，在建规模 710 公里；同一时期，陕西省交通运输系统“十四五”累计完成综合交通投资 4142 亿元，公路、铁路、机场等综合交通基础设施建设保持高位推进。这意味着陕西已经形成了覆盖关中平原城市群、陕北能源运输走廊以及陕南山岭重丘区的高速公路网络结构，既有高流量、强通勤、强物流的平原通道，也有桥隧比高、气象条件复杂、灾害耦合风险更突出的山地通道。

从日常运行特征看，陕西高速公路路网呈现出“枢纽走廊高流量 + 山区通道高风险”并存的双重属性。陕西省交通运输厅持续发布的路网运行周报显示，G3002 西安绕城高速、S1 新机场线、G3021 临潼—兴平高速、G30 连霍高速、S20 西安外环高速等线路长期位居全省交通量排名前列，这类路段监控视频资源丰富、事件影响范围大，对视频感知系统的稳定性、实时性和误报控制能力要求较高。与此同时，陕南地区及秦巴山地带又承担着跨省出行、山区旅游、能源和民生物资运输等任务，路段环境与平原高速存在显著差异，恶劣天气状态下的风险感知更依赖精细化监测。

从安全管理现状看，陕西交通部门已经将“监测预警前置、路网协同联动、应急力量预置”作为高速公路运行管理的重要方向。2025 年 12 月 15 日，陕西省交通运输厅与重庆市交通委员会联合组织了陕渝省际高速公路低温雨雪冰冻天气暨隧道突发事件应急演练，演练场景覆盖暴雪预警、隧道交通事故、车辆受困、交通中断、监测预警、会商研判和交通管控等多个环节。这类省际山地通道演练说明，当前陕西交通安全管理已不仅关注“看见事件”，更关注“提前感知风险、提前组织联动、提前优化处置”的全过程管理能力。

此外，陕西在公路运行监测基础设施上也已具备较好的信息化底座。陕西省交通运输厅公开的 2025 年预算资料中已列出“陕西省高速公路网运行监测能力提升”项目，陕西交通控股集团 2025 年高速公路监测预警系统建设（第一批）招标公告则进一步表明，全省已开始针对 18 家路段分公司的 142 处风险点建设监测预警系统。这些信息说明，本文项目若与陕西省交通运输部门现状相结合，不应被理解为“从零开始替代现有系统”，而更适合定位为：在既有监控视频、路网监测、收费中心和“一路多方”协同体系基础上，补充一层面向低能见度场景的视频智能感知能力，为高速公路团雾识别、风险分级、限速诱导和联动会商提供更细粒度的视觉证据。

\section{研究问题与技术挑战}

针对高速公路团雾监测，本文关注的核心问题可以概括为：如何在固定道路监控视频中，利用可获取的清晰天气交通数据构建可训练样本，并设计一个能够同时输出车辆检测、雾类型分类和雾浓度估计结果的统一模型，从而形成可部署的视频监测原型系统。

围绕上述问题，实际研究面临以下几类关键挑战。

第一，真实团雾样本稀缺且标注成本高。真正发生严重团雾的监控视频出现频率并不高，且容易受到时段、季节、地域和设备条件影响。即便采集到了雾视频，也难以获得精确的每帧雾浓度标签和一致的客观能见度标注。相比之下，已有公开道路监控数据集大多集中于晴天或一般天气条件，难以直接支撑雾浓度回归任务。

第二，雾天气具有明显的空间非均匀性。许多传统去雾或雾检测研究将整幅图像视作均匀退化过程，而高速公路团雾往往具有局部增强、片状分布和随深度变化的非均匀特征。若模型只基于全局平均亮度或简单手工特征，很难稳定地区分无雾、均匀雾和局地团雾。

第三，监控视频与自动驾驶前视相机场景差异明显。许多国际研究集中于车载前视摄像头、合成雾数据集或自动驾驶多传感器融合场景，而高速公路监控摄像头通常具有固定机位、大透视变化、远距离小目标多、图像压缩明显、清晰度受设备限制等特点。针对自动驾驶场景训练的模型并不能直接迁移到固定监控画面。

第四，多任务学习中的目标冲突问题突出。车辆检测关注局部目标、边缘和多尺度结构；雾类型识别更关注整帧全局统计与纹理退化；雾浓度回归需要在视觉退化特征与物理模型之间建立稳定联系。这三类任务在特征偏好和损失尺度上并不一致，如果缺少合理的共享结构与损失设计，模型容易出现其中某一任务主导训练而导致其他任务性能下降。

第五，工程部署要求算法链路完整且具备可解释性。单一的分类结果难以满足道路管控需求；单一的目标检测结果又无法说明天气风险；单一的去雾结果虽改善视觉观感，但并不天然等价于风险量化。因此，实际系统更需要“可见结果 + 可解释指标 + 可扩展部署路径”的综合解决方案。

\section{研究目标与主要内容}

围绕上述问题，本文设定以下研究目标：

\begin{enumerate}[leftmargin=2.5em]
\item 构建一套面向高速公路团雾监测的数据处理链路，能够在缺少大规模真实团雾标注样本的前提下，借助现有晴天交通数据生成适合多任务训练的数据流。
\item 设计一种统一多任务模型，实现车辆检测、雾类型识别和雾浓度回归的联合建模。
\item 实现从训练、验证、视频推理到导出部署的完整工程闭环，并在真实下载视频片段上完成系统级验证。
\item 为后续毕业论文答辩、论文扩写和实验补充提供可复现的工程原型、技术文本与方法论基础。
\end{enumerate}

为达成以上目标，本文主要工作包括：

\begin{enumerate}[leftmargin=2.5em]
\item 分析高速公路团雾监测的业务需求与视觉任务特点，明确“检测 + 分类 + 回归”三任务联合建模的必要性；
\item 基于 UA-DETRAC 数据集构建原始数据索引与序列级训练/验证划分，并使用 \code{MiDaS_small} 预计算深度缓存；
\item 设计基于大气散射模型的在线造雾模块，在 GPU 侧动态生成 clear、uniform、patchy 三类天气样本；
\item 构建以 YOLO 检测骨干为共享主干的统一多任务模型，在高层语义特征上附加雾分类头与 beta 回归头；
\item 实现带有 checkpoint 续训、自动补算深度缓存、混合精度训练、QAT 路径和视频演示推理的训练与部署代码；
\item 设计实验方案，对数据、方法、系统功能和工程可行性进行综合分析，并在后续训练结果完善后补充完整量化结论。
\end{enumerate}

\section{本文创新点}

结合项目目标与当前实现，本文的创新点主要体现在以下几个方面。

第一，提出了一种面向固定交通监控视频的“深度缓存 + 在线物理造雾”数据构建策略。该方法不需要预先生成海量雾图，也不强依赖大规模真实雾天气标注，而是借助深度估计和大气散射模型，在训练阶段动态生成不同类型和强度的雾样本，从而降低数据获取成本并提升样本多样性。

第二，设计了一种面向高速公路团雾监测的统一多任务模型。与传统仅进行车辆检测或仅进行天气分类的方法不同，本文将车辆检测、雾类型分类和雾浓度回归整合到同一网络中，以共享高层语义特征实现三任务协同学习，增强了模型的整体表达能力和系统一体化程度。

第三，从工程可落地角度提出了“统一模型主链 + 混合推理备份链”的系统实现方案。当统一模型在训练早期对车辆检测的鲁棒性尚未充分提升时，可通过 YOLO11n 检测器与雾估计头组合形成混合推理路径，保证系统演示与业务验证不断链。这种设计体现了从纯学术模型走向工程系统时对稳定性与可交付性的综合考虑。

第四，建立了较为完整的训练与部署工具链。本文不仅完成了网络结构设计，还实现了数据审计、深度缓存预计算、训练前 smoke test、视频推理、ONNX 导出以及 QAT/INT8 路径，为后续边缘部署和性能优化留下了明确扩展接口。

\section{论文结构安排}

本文全文共分为九章。第一章介绍研究背景、问题定义、研究目标和创新点；第二章综述国内外在团雾监测、能见度识别、图像去雾和恶劣天气目标检测方面的研究进展；第三章介绍大气散射模型、单目深度估计、多任务学习和目标检测等理论基础；第四章从工程视角给出系统需求分析与总体技术方案；第五章详细说明数据集构建、深度缓存和在线造雾方法；第六章重点阐述统一多任务模型结构和训练优化策略；第七章介绍视频推理、可视化展示与部署导出实现；第八章给出实验设计、当前结果和后续补充计划；第九章总结全文并展望未来工作。
